% Recovered from direct-forced-intersection.md, lines 466--516, and
% navier-stokes-pressure-material-investigation.md, lines 141--212,
% 553--763, and 821--868.  The main paper uses B_n for full transport.
% This fragment preserves that notation; projected transport is \Pp B_n.
% All norms in time are over (0,T).  Insert after the basic rectangle,
% product, pressure-projection, and acceleration identities are available.

\section{General bounds for production, pressure, and acceleration}\label{r:sec:recovered-rectangle}
\label{r:sec:recovered-rectangle-bounds}

The selected rectangles control more than the critical energy and the
pointwise momentum terms.  Each temporal row has its own nonlinear
production and source work.  Retaining the full binomial row lets us bound
these quantities together, and then use the two pressure projections to
resolve how transport contributes to material acceleration.

\subsection{The production and work of all temporal rows}

Fix nonnegative integers \(N,m\).  The pressure work in row \(r\) vanishes
after pairing with the solenoidal field \(V_r\).  The product estimate
\eqref{r:eq:transport-product} gives
\begin{equation}
 |\mathcal P_{r,m}|
 \le C_m\norm{V_r}_{H^m}
 \sum_{a=0}^r\binom ra
 \norm{V_a}_{H^{m+2}}\norm{V_{r-a}}_{H^{m+2}},
 \qquad C_m=2^{m+1}c_\infty.
 \label{r:eq:all-row-production-point}
\end{equation}
The two factors in every summand are retained: temporal differentiation
does not preserve every cancellation of the zeroth row.

Write
\[
 D_{N,m}:=\sum_{r=0}^N\norm{V_r(0)}_{H^m}^2.
\]
Integrating the adjacent trace identity from time zero and using
\(2ab\le a^2+b^2\) gives
\begin{equation}
 \max_{r\le N}\norm{V_r}_{L^\infty H^m}
 \le\sqrt{D_{N,m}+R_{N,m}}.
 \label{r:eq:all-row-initial-trace}
\end{equation}
For \(x_a=\norm{V_a}_{L^2H^{m+2}}\), time Cauchy--Schwarz and the
binomial symmetry give
\[
 \sum_{a=0}^r\binom ra x_ax_{r-a}
 \le 2^r\sum_{a=0}^N x_a^2
 \le 2^rR_{N,m+1}.
\]
We can now integrate \eqref{r:eq:all-row-production-point} and sum the full nonlinear contribution through row \(N\):
\begin{equation}
 \sum_{r=0}^N\norm{\mathcal P_{r,m}}_{L^1}
 \le C_m(2^{N+1}-1)
 \sqrt{D_{N,m}+R_{N,m}}\,R_{N,m+1}.
 \label{r:eq:all-row-production-integral}
\end{equation}
This bounds the absolute production in every selected row, including
each signed prefix integral.

The source work admits a bound with one adjacent spatial derivative
placed on each side of the pairing.  Since
\[
 W_{r,m}
 =\operatorname{Re}\ip{\Lambda^{m+1}V_r}{\Lambda^{m-1}F_r},
\]
define
\[
 \mathcal F_{N,m-1}^{\rm work}
 :=\left(\sum_{r=0}^N
 \norm{\Lambda^{m-1}F_r}_{L^2L^2}^2\right)^{1/2}.
\]
Cauchy--Schwarz first in time and then over the finite rows yields
\begin{equation}
 \sum_{r=0}^N\norm{W_{r,m}}_{L^1}
 \le\mathcal F_{N,m-1}^{\rm work}\sqrt{R_{N,m}}.
 \label{r:eq:all-row-force-work}
\end{equation}
At \(m=0\), the source coordinate is \(\dot H^{-1}\).  The assumed
spatial decay supplies its low frequencies.  For each fixed temporal
order \(\ell\), differentiation of the pairing gives the further bound
\begin{equation}
 \begin{aligned}
 \norm{\partial_t^\ell W_{r,m}}_{L^1}
 &\le\sum_{a=0}^{\ell}\binom\ell a
 \norm{\Lambda^{m+1}V_{r+a}}_{L^2L^2}
 \norm{\Lambda^{m-1}F_{r+\ell-a}}_{L^2L^2}.
 \end{aligned}
 \label{r:eq:all-row-force-work-derivatives}
\end{equation}
Every factor is a finite coordinate of the velocity or the prescribed
source.  The constants may depend on \(N,m,\ell\); the calculation
requires no bound uniform over all derivative orders.

\subsection{Projected transport from products and adjacent rows}

The pressure projection also gives direct integral estimates for the
nonlinear field, before it is paired with velocity.  Put
\[
 X_{a,j}:=\norm{V_a}_{L^2H^j}.
\]
For every nonnegative integer \(M\), choose the finite product constant
\(C_M^{\rm alg}\) in
\[
 \norm{v\otimes w}_{H^M}
 \le C_M^{\rm alg}\norm v_{H^{M+1}}\norm w_{H^{M+1}}.
\]
For \(M=0\), this follows from \(H^1\subset L^4\) and H\"older;
at higher integer orders it follows by the finite Leibniz expansion
and the Sobolev product bounds.  Tensor divergence maps \(H^M\) to
\(H^{M-1}\) with norm at most one, and both \(\Pp\) and
\(\mathbb Q=I-\Pp\) are contractions in these Fourier Sobolev norms.
Applying them to the full transport row gives
\begin{equation}
 \begin{aligned}
 \norm{\Pp B_n}_{L^1H^{M-1}},\quad
 \norm{\mathbb Q B_n}_{L^1H^{M-1}}
 &\le C_M^{\rm alg}
 \sum_{a=0}^n\binom na X_{a,M+1}X_{n-a,M+1}\\
 &\le C_M^{\rm alg}2^nR_{N,M},\qquad n\le N.
 \end{aligned}
 \label{r:eq:projected-transport-product-integral}
\end{equation}
Here and below the two quantities separated by a comma each obey the
displayed bound.

Let \(Q_n^u\) denote the normalized nonlinear pressure in row \(n\),
so that \(\nabla Q_n^u=-\mathbb Q B_n\).  Its scalar reconstruction
is a double Riesz transform of the binomial tensor sum.  The multiplier
has norm at most one from the Frobenius tensor norm to the scalar
norm.  The product estimate consequently gives
\begin{equation}
 \begin{aligned}
 \norm{Q_n^u}_{L^1H^M}
 +\norm{\nabla Q_n^u}_{L^1H^{M-1}}
 &\le 2C_M^{\rm alg}
 \sum_{a=0}^n\binom na X_{a,M+1}X_{n-a,M+1}\\
 &\le 2C_M^{\rm alg}2^nR_{N,M}.
 \end{aligned}
 \label{r:eq:nonlinear-pressure-product-integral}
\end{equation}
The total pressure gradient includes \(\mathbb QF_n\) as well.

Using the next temporal coordinate gives a stronger time norm for the
projected transport.  The momentum row itself reads
\[
 \Pp B_n-\Pp F_n=\nu\Delta V_n-V_{n+1}.
\]
Its two terms on the right belong to the adjacent rectangle, so
\[
 \norm{\Pp B_n-\Pp F_n}_{L^2H^m}
 \le\nu X_{n,m+2}+X_{n+1,m}.
\]
Scalar Cauchy--Schwarz with coefficients \((\nu,1)\), followed by
summation over \(n\), gives the explicit complete-row bound
\begin{equation}
 \left(\sum_{n=0}^N
 \norm{\Pp B_n-\Pp F_n}_{L^2H^m}^2\right)^{1/2}
 \le\sqrt{\nu^2+1}\,R_{N,m+1}^{1/2}.
 \label{r:eq:projected-transport-adjacent}
\end{equation}
The right side retains the next temporal row.  This is the direct
bound of the nonlinear field participating in the equation along the
selected history.

The spatial recovery in that equation is also explicit.  The Fourier
inequality \((1+|\xi|^2)^2\le2(1+|\xi|^4)\) implies
\begin{equation}
 \norm{V_n}_{H^{m+2}}
 \le\sqrt2\left[
 \norm{V_n}_{H^m}
 +\nu^{-1}\norm{V_{n+1}+\Pp B_n-\Pp F_n}_{H^m}\right].
 \label{r:eq:adjacent-spatial-recovery}
\end{equation}
Both the temporal derivative and the complete projected nonlinear
response remain in this recovery of two spatial derivatives.

\subsection{Higher pressure rows and their time norms}

The pressure gain in \eqref{r:eq:pressure-critical-gain} extends to every
temporal row.  For two solenoidal fields, expanding the divergence gives
\[
 \operatorname{div}((v\cdot\nabla)w)
 =\partial_i v_j\,\partial_j w_i.
\]
This expression is symmetric under exchanging \(v,w\).  The gradient
projection has the exact norm relation
\[
 \norm{\mathbb Q((v\cdot\nabla)w)}_{\dot H^{1/2}}
 =\norm{\partial_i v_j\,\partial_j w_i}_{\dot H^{-1/2}}.
\]
With \(c_*=(\sqrt2\pi)^{-1}\), the Sobolev inequality and its dual
bound the right side by
\(c_*\norm v_{\dot H^{3/2}}\norm w_{\dot H^{3/2}}\).
Apply this to every binomial summand:
\begin{equation}
 \norm{\nabla Q_n^u}_{\dot H^{1/2}}
 \le c_*\sum_{a=0}^n\binom na
 \norm{V_a}_{\dot H^{3/2}}\norm{V_{n-a}}_{\dot H^{3/2}}.
 \label{r:eq:pressure-all-temporal-endpoint}
\end{equation}
The symmetry also gives
\begin{equation}
 \begin{aligned}
 \norm{\mathbb Q((v\cdot\nabla)w)}_{\dot H^{-1/2}}
 \le c_*\min\bigl\{&
 \norm v_{\dot H^{1/2}}\norm w_{\dot H^{3/2}},\\
 &\norm w_{\dot H^{1/2}}\norm v_{\dot H^{3/2}}\bigr\}.
 \end{aligned}
 \label{r:eq:pressure-factor-exchange}
\end{equation}
Either assignment of the two heights is available for the pressure
component because its divergence retains a derivative on both factors.

For the zeroth row, put
\[
 h(t)=\norm{u(t)}_{\dot H^{1/2}},\qquad
 d(t)=\norm{u(t)}_{\dot H^{3/2}},\qquad
 H_*:=\norm h_{L^\infty},\quad D_*:=\norm d_{L^2}.
\]
The bounds \(\norm{B_0}_{\dot H^{-1/2}}\le c_*hd\) and
\(\norm{\mathbb QB_0}_{\dot H^{1/2}}\le c_*d^2\) give
\begin{equation}
 \begin{aligned}
 \norm{\Pp B_0}_{L^2\dot H^{-1/2}},\quad
 \norm{\mathbb QB_0}_{L^2\dot H^{-1/2}}
 &\le c_*H_*D_*,\\
 \norm{\mathbb QB_0}_{L^1\dot H^{1/2}}
 &\le c_*D_*^2.
 \end{aligned}
 \label{r:eq:pressure-time-endpoints}
\end{equation}
Interpolating the two spatial norms at each time gives
\[
 \norm{\mathbb QB_0(t)}_2
 \le c_*h(t)^{1/2}d(t)^{3/2}.
\]
Raising this to the power \(4/3\) and integrating yields
\begin{equation}
 \norm{\mathbb QB_0}_{L^{4/3}L^2}
 \le c_*H_*^{1/2}D_*^{3/2}.
 \label{r:eq:pressure-intermediate-time}
\end{equation}
The critical rectangle supplies \(H_*\le\sqrt{D+R}\) and
\(D_*\le\sqrt R\), so the last right side is
\(c_*(D+R)^{1/4}R^{3/4}\).  A Sobolev and H\"older estimate also
gives the \(L^{4/3}L^2\) conclusion for full transport.  The additional
pressure regularity is its \(L^1\dot H^{1/2}\) endpoint; total forced
pressure still includes the prescribed gradient source.

\subsection{Material acceleration from the first rectangle}

The orthogonal acceleration identity can be evaluated in \(L^2_x\)
using \(R_{0,1}\).  Set
\[
 \mathsf A:=\norm u_{L^2H^2},\qquad
 \mathsf B:=\norm{u_t}_{L^2L^2},\qquad
 \mathsf Z:=\left(\norm{u_0}_{H^1}^2
                   +2\mathsf A\mathsf B\right)^{1/2}.
\]
The adjacent trace gives \(\norm u_{L^\infty H^1}\le\mathsf Z\).
Use the homogeneous Sobolev constants
\[
 S_6=\frac{2^{2/3}}{\sqrt3\pi^{2/3}},\qquad
 S_3=2^{-1/6}\pi^{-1/3}.
\]
H\"older in space and then in time gives
\begin{equation}
 \begin{aligned}
 \norm{B_0}_{L^2L^2}
 &\le\norm u_{L^\infty L^6}\norm{\nabla u}_{L^2L^3}\\
 &\le S_6S_3\mathsf Z\mathsf A
 =\frac1\pi\sqrt{\frac23}\,\mathsf Z\mathsf A.
 \end{aligned}
 \label{r:eq:first-rectangle-transport-L2}
\end{equation}
In unforced flow, viscosity supplies the solenoidal part of acceleration
and the pressure response supplies its gradient part.  Their squared
norms add by \eqref{r:eq:acceleration-orthogonal}.  We obtain
\begin{equation}
 \norm{D_tu}_{L^2L^2}
 \le\mathsf A\sqrt{\nu^2+\frac{2\mathsf Z^2}{3\pi^2}},
 \qquad f=0.
 \label{r:eq:first-rectangle-acceleration-L2}
\end{equation}
For the forced equation, retain the prescribed solenoidal source in
\eqref{r:eq:acceleration-orthogonal}. The resulting bound at this first
rectangle is
\begin{equation}
 \norm{D_tu}_{L^2L^2}^2
 \le\left(\nu\mathsf A+\norm{\Pp f}_{L^2L^2}\right)^2
 +\frac{2}{3\pi^2}\mathsf Z^2\mathsf A^2.
 \label{r:eq:first-rectangle-forced-acceleration-L2}
\end{equation}
The triangle inequality bounds the solenoidal component, while the
preceding transport estimate bounds the orthogonal gradient component.
The adjacent row also gives
\(\norm{\Pp B_0}_{L^2L^2}\le\nu\mathsf A+\mathsf B\) when
\(f=0\).  These bounds use the first rectangle, while the earlier
pointwise acceleration bound selects the stronger \(R_{1,2}\).

At critical negative-half order, the orthogonal identity gives a
particularly precise unforced coefficient:
\begin{equation}
 \nu d(t)\le\norm{D_tu(t)}_{\dot H^{-1/2}}
 \le\sqrt{\nu^2+\frac{h(t)^2}{2\pi^2}}\,d(t),
 \qquad f=0.
 \label{r:eq:critical-acceleration-two-sided}
\end{equation}
The lower bound retains the solenoidal viscous component, and the upper
bound uses \(\norm{\mathbb QB_0}_{\dot H^{-1/2}}\le c_*hd\).
Integration gives
\[
 \norm{D_tu}_{L^2\dot H^{-1/2}}
 \le\sqrt{\nu^2+\frac{H_*^2}{2\pi^2}}\,D_*,\qquad f=0.
\]
For general forcing, the solenoidal component is
\(\nu\Delta u+\Pp f\); retaining it gives
\begin{equation}
 \norm{D_tu}_{L^2\dot H^{-1/2}}^2
 \le\left(\nu D_*+\norm{\Pp f}_{L^2\dot H^{-1/2}}\right)^2
       +c_*^2H_*^2D_*^2.
 \label{r:eq:critical-acceleration-general-coordinates}
\end{equation}
Its substitution from \(R_{0,1}\) is
\eqref{r:eq:critical-material-bounds}.

\subsection{Acceleration at a higher adjacent height}

To place the full acceleration in \(L^2H^{1/2}\), select the stronger
adjacent coordinates
\begin{equation}
 u\in L^2H^{5/2},\qquad u_t\in L^2H^{1/2}.
 \label{r:eq:stronger-acceleration-coordinates}
\end{equation}
Higher integer rectangles supply both by Sobolev inclusion.  The
inhomogeneous pairing
\(2\ip{(1-\Delta)^{5/4}u}{(1-\Delta)^{1/4}u_t}\)
gives the continuous \(H^{3/2}\) trace and
\begin{equation}
 \sup_{0\le t\le T}\norm{u(t)}_{H^{3/2}}^2
 \le\norm{u_0}_{H^{3/2}}^2
 +2\norm u_{L^2H^{5/2}}\norm{u_t}_{L^2H^{1/2}}.
 \label{r:eq:stronger-acceleration-trace}
\end{equation}
For the pressure component, the negative-half estimate controls
frequencies below one and the positive-half estimate controls those
above one.  Together they give a continuous quadratic map from
\(H^{3/2}\) into \(H^{1/2}\).  Thus
\(\mathbb QB_0\in C([0,T];H^{1/2})\), with
\begin{equation}
 \norm{\mathbb QB_0}_{L^2H^{1/2}}
 \le C\norm u_{L^\infty H^{3/2}}
         \norm u_{L^2H^{3/2}}.
 \label{r:eq:stronger-pressure-time}
\end{equation}
If \(\Pp f\in L^2H^{1/2}\), all components of
\(a=\nu\Delta u+\Pp f+\mathbb QB_0\) now belong to
\(L^2H^{1/2}\).  The projected momentum row separately yields
\begin{equation}
 \norm{\Pp B_0}_{L^2H^{1/2}}
 \le\nu\norm u_{L^2H^{5/2}}
 +\norm{u_t}_{L^2H^{1/2}}
 +\norm{\Pp f}_{L^2H^{1/2}}.
 \label{r:eq:stronger-projected-transport}
\end{equation}
The \(L^2H^{5/2}\) bound controls the viscous term in
\(L^2H^{1/2}\), while the adjacent temporal bound yields
the continuous \(H^{3/2}\) trace.

\subsection{Unforced action and spatial recovery}

For \(f=0\), squaring the projected equation in
\(\dot H^{-1/2}\) retains an exact relation between the temporal
and viscous actions.  Since
\[
 \Pp B_0=\nu\Delta u-u_t,\qquad
 \ip{\Delta u}{u_t}_{\dot H^{-1/2}}
 =-\frac12\frac{\dd}{\dd t}h^2,
\]
we have
\begin{equation}
 \nu\frac{\dd}{\dd t}h^2+\nu^2d^2
 +\norm{u_t}_{\dot H^{-1/2}}^2
 =\norm{\Pp B_0}_{\dot H^{-1/2}}^2,
 \qquad f=0.
 \label{r:eq:unforced-critical-action}
\end{equation}
Add \(\norm{\mathbb QB_0}_{\dot H^{-1/2}}^2\) to both sides.
Orthogonality of the two transport projections and the acceleration
identity give
\begin{equation}
 \nu\frac{\dd}{\dd t}h^2
 +\norm{u_t}_{\dot H^{-1/2}}^2
 +\norm{D_tu}_{\dot H^{-1/2}}^2
 =\norm{B_0}_{\dot H^{-1/2}}^2,
 \qquad f=0.
 \label{r:eq:unforced-material-critical-action}
\end{equation}
These identities retain the complete nonlinear action on the right.
Adding the pressure square also adds its contribution there; the
operation supplies no additional sign for critical production.

The solenoidal acceleration determines the viscous spatial derivative
when \(\nu>0\).  In unforced flow,
\(\Pp a=\nu\Delta u\), so the Fourier estimate used in
\eqref{r:eq:adjacent-spatial-recovery} gives
\begin{equation}
 \norm u_{H^{m+2}}
 \le\sqrt2\left(\norm u_{H^m}
                  +\nu^{-1}\norm a_{H^m}\right),
 \qquad f=0.
 \label{r:eq:unforced-spatial-recovery-from-acceleration}
\end{equation}
This last relation recovers the spatial coordinate from the full
acceleration at that height, with the lower velocity norm retained.
